

\begin{document}

\section{Introduction}
T cell receptors (TCRs) play a central role in adaptive immunity by recognizing antigenic peptides presented by major histocompatibility complex (MHC) molecules. Characterizing TCR–peptide–MHC (pMHC) interactions is therefore fundamental to understanding antigen-specific immune responses and has important applications in infectious disease, autoimmunity, cancer immunotherapy, and vaccine development [REF: general TCR–pMHC recognition/immunology]. Experimentally, TCR specificity can be characterized using pMHC tetramer staining, functional assays, and more recent high-throughput approaches capable of screening large numbers of TCRs and antigens [REF: pMHC tetramer methods; high-throughput TCR specificity mapping]. Although these approaches provide direct experimental evidence of TCR recognition, exhaustive experimental characterization is difficult because of the enormous combinatorial space of TCR and antigen sequences.

Computational prediction has consequently emerged as an important complementary approach for studying TCR–antigen recognition. Over the past decade, numerous computational models have been developed to infer TCR specificity from sequence information, ranging from similarity-based and conventional machine-learning approaches to deep neural networks and representation-learning methods [REF: TCR prediction review; representative early models]. Representative methods include ERGO/ERGO-II, NetTCR, TITAN, PanPep, SCEPTR, EPACT, and related approaches [REF: original model papers]. In principle, the utility of these models extends beyond predicting individual TCR–epitope interactions. A reliable computational model could be used to explore antigen sequence space at a scale that is difficult to achieve experimentally, predict how sequence variation affects existing T-cell recognition, and potentially identify antigen mutations associated with immune escape or reduced vaccine-induced recognition [REF: antigenic variation/T-cell escape; computational mutation prediction if available].

Despite this potential, current TCR–epitope prediction is predominantly formulated and evaluated as a binary pairwise classification problem, in which a TCR and peptide are treated as an independent pair and the model predicts whether they bind or do not bind. Recent large-scale benchmarking efforts have substantially improved our understanding of the limitations of this formulation. The IMMREP23 community benchmark demonstrated a pronounced reduction in predictive performance for previously unseen epitopes [REF: Nielsen/Eugster et al., IMMREP23]. More recently, Lu et al. systematically evaluated a broad range of TCR–epitope prediction methods and similarly showed that performance depends strongly on training-data composition, negative-sample construction, sequence similarity, and whether test epitopes are represented in training data [REF: Lu et al., Nature Methods]. These studies establish that generalization remains a major challenge for current TCR prediction models.

However, successful discrimination between binding and non-binding TCR–epitope pairs addresses only one aspect of the biological recognition problem. Antigen sequences evolve and vary, and even a single amino-acid substitution can alter TCR recognition. Depending on the location and identity of the substitution and the particular TCR involved, a mutation may preserve recognition, weaken it, or result in loss of detectable binding [REF: TCR cross-reactivity/peptide mutation literature]. Such mutation-dependent recognition is directly relevant to viral evolution, immune escape, vaccine protection, and TCR cross-reactivity [REF: viral T-cell escape; vaccine escape/cross-reactivity]. Importantly, mutation effects are structured rather than uniform: some peptide positions tolerate broad sequence variation, whereas others are strongly constrained, and the effects at many positions are TCR-dependent [REF: peptide scanning/TCR recognition; Malone et al.]. Thus, an important question remains largely unexplored by conventional benchmarks: do computational TCR–epitope prediction models reproduce experimentally observed changes in TCR recognition caused by antigen mutations?

Answering this question also requires careful attention to the reliability of the computational evaluation itself. TCR datasets collected from nominally independent studies can contain identical or highly similar TCRs and overlapping epitopes, allowing training–test leakage to inflate apparent prediction performance [REF: IMMREP23; Lu et al.; other leakage studies]. This problem becomes particularly important when evaluating antigen mutations because the test peptides are intentionally highly similar to one another and to the wild-type antigen. In addition, many TCR prediction datasets lack experimentally established non-binding interactions and therefore rely on synthetic negatives constructed by randomly mismatching known TCRs and peptides [REF: negative sampling studies]. Such pairs cannot necessarily be regarded as biological non-binders because unmeasured cross-reactivity may exist. A rigorous evaluation of mutation effects should therefore combine experimentally determined recognition measurements with explicit control of each model's prior exposure to the evaluated TCRs and antigens.

A recent experimental study by Malone et al. provides an opportunity to address these limitations [REF: Malone et al., Nature Communications, 2026]. Using a TCR fingerprinting approach, the study systematically characterized recognition of the immunodominant SARS-CoV-2 epitope YLQPRTFLL (YLQ) and its single-amino-acid variants. Twenty-one TCRs were evaluated against a panel of 172 peptide variants spanning substitutions across the nine positions of YLQ. TCR-specific pMHC recognition was measured by staining TCR-transduced J76 cells and normalized against TCR-negative J76 cells, producing quantitative measurements expressed as log
2
	​

(signal MFI/background MFI). The experimental study revealed highly structured mutation-response profiles, including broadly tolerated peptide positions, strongly restrictive positions, and positions exhibiting substantial TCR-specific variation [REF: Malone et al.]. These measurements therefore provide substantially richer information than collections of unrelated binary TCR–epitope pairs: for each TCR, they define an experimentally measured recognition fingerprint across a controlled antigen mutation landscape.

In this study, we transform these experimental measurements into a mutation-aware computational benchmark for evaluating TCR–epitope prediction models. The benchmark preserves the systematic relationship among the wild-type antigen and its single-amino-acid variants, experimentally determined binding states, quantitative pMHC-binding measurements, mutation positions, and individual TCR identities. We further develop a model-independent, leakage-free evaluation pipeline that compares computational predictions with experimental observations at multiple resolutions. In addition to conventional interaction-level performance, the framework evaluates whether a model correctly predicts the retention or loss of recognition for individual mutations, recovers position-specific mutation effects, and reproduces TCR-specific recognition fingerprints. By explicitly accounting for each model's training data, the framework is designed to distinguish genuine mutation-prediction capability from performance attributable to prior exposure to overlapping TCRs or antigens.

We selectively apply this framework to eight widely used TCR–epitope prediction models representing different computational approaches [REF: eight model papers]. [INSERT PRINCIPAL RESULTS AFTER COMPLETING THE MUTATION-WISE AND POSITION-WISE ANALYSES; e.g., “We find that current models show limited ability to reproduce experimentally observed mutation-induced changes in TCR recognition, with substantial variation across antigen positions and TCRs. Although some models capture selected coarse positional constraints, their predictions generally fail to reproduce the experimentally observed TCR-specific recognition fingerprints.”] These results suggest that performance on conventional TCR–epitope binding benchmarks does not necessarily imply an ability to predict the biological consequences of antigen sequence variation.

The main contributions of this study are threefold. 
\begin{itemize}
    \item First, we introduce a new mutation-aware computational benchmark derived from systematic experimental antigen perturbations, enabling direct evaluation of how TCR prediction models respond to single-amino-acid antigen mutations. 
    \item Second, we develop a leakage-free, model-independent evaluation pipeline for quantifying agreement between computational predictions and experimentally observed mutation effects at the individual-mutation, antigen-position, and TCR-specific recognition-profile levels. 
    \item Third, through a systematic evaluation of eight representative TCR–epitope prediction models, we characterize [INSERT FINAL PRINCIPAL FINDING] and identify capabilities and limitations of current models that are not revealed by conventional aggregate binding-prediction benchmarks. 
\end{itemize}
Together, the benchmark and evaluation framework extend computational assessment of TCR specificity from predicting whether a TCR recognizes an antigen toward determining whether models can capture how antigen sequence variation changes TCR recognition.


\section{Materials and Methods}

\subsection{Overview of the Mutation-Aware Benchmarking Framework}

This study does not propose a new T cell receptor (TCR)--epitope 
prediction model. Instead, we develop a model-independent benchmarking 
framework to evaluate whether existing computational models reproduce 
experimentally observed effects of antigen mutations on TCR recognition. 
The overall workflow is illustrated in Fig.~\ref{fig:framework}.

The framework consists of five major steps: 
(1) construction of a computational benchmark from experimentally measured TCR recognition across systematic single-amino-acid antigen substitutions; 
(2) collection and examination of the training data used by individual prediction models to identify potential training--test overlap; 
(3) inference of TCR--epitope binding scores using selected prediction models; 
(4) evaluation of mutation effects at the individual-interaction, individual-mutation, antigen-position, and TCR-specific levels; and 
(5) comparison of computational predictions with both binary and quantitative 
experimental measurements. This design enables evaluation of whether a prediction model captures the effects of antigen sequence variation rather 
than only its aggregate ability to distinguish binding from non-binding 
TCR--epitope pairs.


\subsection{Experimental Fingerprinting Dataset}

\subsubsection{TCRs and antigen mutation panel}
The experimental data used in this study were generated using the TCR fingerprinting approach reported by Malone et al.~\cite{malone2025fingerprinting}. 
The study characterized recognition of the immunodominant SARS-CoV-2 epitope YLQPRTFLL (YLQ) and a systematic panel of single-amino-acid variants. 
A total of 21 TCRs were evaluated against 172 peptide variants spanning substitutions across the nine amino-acid positions of YLQ.

Unlike conventional TCR--epitope datasets assembled from unrelated positive and negative interactions, the Fingerprinting data represent controlled perturbations of a common antigen sequence. Each mutant therefore differs from the wild-type antigen by a defined single-amino-acid substitution, allowing the effect of antigen sequence variation on TCR recognition to be examined systematically. The resulting data provide a structured antigen mutation landscape for each evaluated TCR.

% INSERT DETAILS IF NEEDED:
% - source and characteristics of the 21 TCRs
% - number of substitutions represented at each position
% - excluded peptide variants and reasons for exclusion


\subsubsection{Experimental pMHC-binding measurements}
TCR-specific peptide--MHC (pMHC) recognition was measured by staining TCR-transduced J76 cells with individual pMHC tetramers. Background staining 
was measured using TCR-negative J76 cells. For TCR $i$ and peptide variant $j$, the quantitative experimental binding signal was represented as

\begin{equation}
B_{ij}=
\log_2
\left(
\frac{\mathrm{MFI}_{\mathrm{TCR}^{+},ij}}
     {\mathrm{MFI}_{\mathrm{TCR}^{-},j}}
\right),
\label{eq:binding_signal}
\end{equation}

where $\mathrm{MFI}_{\mathrm{TCR}^{+},ij}$ denotes the pMHC staining 
intensity for TCR-transduced cells and $\mathrm{MFI}_{\mathrm{TCR}^{-},j}$ 
denotes the corresponding background staining intensity measured using TCR-negative cells on peptide variant $j$.


To minimize the possibility that apparent differences in TCR recognition resulted from insufficient peptide--MHC binding or unstable pMHC complexes, the original study evaluated the HLA-A*02:01 binding capacity of the peptide
variants. Peptides were first assessed using NetMHCpan 4.1 \cite{reynisson20}, with predicted weak and strong MHC binders defined by \%Rank $<2$ and $<0.5$, respectively. Peptide--MHC stability was further evaluated experimentally using a T2 stabilization assay, and pMHC species failing the experimentally defined
stability criterion were excluded from downstream analysis \cite{malone2025fingerprinting}. These quality-control procedures help separate mutation-induced changes in TCR recognition from changes caused
primarily by impaired peptide--MHC binding.

For the retained pMHC species, TCR recognition was quantified as the $\log_2$ fold-change in pMHC staining MFI relative to TCR-negative J76 background. Interactions with log$_2$ fold-change $>1$ were classified as binders, corresponding to a staining signal greater than two-fold above
background.
%Following the experimental criteria established by Malone et al.~\cite{malone2025fingerprinting}, the continuous measurements were also converted into binary binding labels. A TCR--pMHC interaction was considered a binder when the staining signal exceeded the experimentally defined threshold:

\begin{equation}
y_{ij} =
\begin{cases}
1, & B_{ij} > 1,\\
0, & B_{ij} \leq 1.
\end{cases}
\label{eq:binary_label}
\end{equation}


Thus, the Fingerprinting data provide two complementary experimental endpoints: a binary measure indicating retention or loss of detectable TCR recognition and a continuous measure describing the strength of pMHC 
staining relative to experimental background.


\subsection{Construction of the Fingerprinting Computational Benchmark}

We transformed the experimental Fingerprinting measurements into a structured computational benchmark for mutation-aware evaluation of TCR--epitope prediction models. Each benchmark record corresponds to a 
specific TCR and peptide variant and contains the information required for both model inference and subsequent mutation-aware analysis. 

Conceptually, each benchmark interaction can be represented as

\begin{equation}
D_{ij}
=
\left(
T_i,P_j, k_j, a_j^{\mathrm{r}}, a_j^{\mathrm{mut}}, B_{ij}, y_{ij}
\right),
\label{eq:benchmark_record}
\end{equation}

where $T_i$ denotes the TCR sequence information, $P_j$ denotes the mutant 
peptide sequence, $k_j$ is the position of the mutation, $a_j^{\mathrm{r}}$ and $a_j^{\mathrm{mut}}$ denote the reference and the
substituted amino acids (the reference epitope is YLQPRTFLL), respectively, $B_{ij}$ is the quantitative 
experimental binding measurement, and $y_{ij}$ is the corresponding binary 
binding label.

The benchmark preserves the relationship among peptide variants derived from the same reference antigen rather than treating TCR--peptide interactions as 
independent observations. This organization enables evaluation at multiple levels, including individual TCR--peptide interactions, individual antigen 
mutations, peptide positions, and complete TCR-specific mutation profiles.



\subsection{Evaluated TCR--Epitope Prediction Models}
We selected eight representative TCR--epitope prediction models spanning 
different computational architectures and training strategies:
ERGO, ERGO-II, NetTCR-2.0, NetTCR-2.2, TITAN, PanPep, SCEPTR, and EPACT
\cite{ergo,ergo2,nettcr20,nettcr22,titan,panpep,sceptr,epact}.
The models were selected to represent widely used approaches to 
sequence-based TCR specificity prediction while providing sufficient 
diversity in model architecture, TCR representation, and training data.

For each model, we recorded the required TCR input (e.g., CDR3$\beta$ or 
paired $\alpha\beta$ chains), peptide and MHC requirements, original training 
data, model output, and recommended decision threshold. Whenever possible, 
publicly released pretrained models and the inference procedures specified 
by the original authors were used.

\begin{table*}[t]
\centering
\caption{TCR--epitope prediction models evaluated in this study.}
\label{tab:models}
\begin{tabular}{llllll}
\hline
Model & TCR input & Peptide input & Training data & Output & Threshold \\
\hline
ERGO       & [TBD] & [TBD] & [TBD] & [TBD] & [TBD] \\
ERGO-II    & [TBD] & [TBD] & [TBD] & [TBD] & [TBD] \\
NetTCR-2.0 & [TBD] & [TBD] & [TBD] & [TBD] & [TBD] \\
NetTCR-2.2 & [TBD] & [TBD] & [TBD] & [TBD] & [TBD] \\
TITAN      & [TBD] & [TBD] & [TBD] & [TBD] & [TBD] \\
PanPep     & [TBD] & [TBD] & [TBD] & [TBD] & [TBD] \\
SCEPTR     & [TBD] & [TBD] & [TBD] & [TBD] & [TBD] \\
EPACT      & [TBD] & [TBD] & [TBD] & [TBD] & [TBD] \\
\hline
\end{tabular}
\end{table*}


\subsection{Model-Specific Leakage Control}

\subsubsection{Collection of model training data}

To minimize the possibility that benchmark performance resulted from prior exposure to the evaluated TCRs or antigens, we examined the training data 
used by each prediction model. Original training datasets released by the model developers were used whenever available. When the exact training data 
were not publicly released, they were reconstructed according to the dataset 
sources and preprocessing procedures described in the corresponding 
publication.

% INSERT TABLE OR SUPPLEMENTARY TABLE:
% Model | Original training dataset | Source | Availability | Reconstruction


\subsubsection{Identification of training--benchmark overlap}

Training--benchmark overlap was assessed separately for each model because 
the models were developed using different training resources. We examined 
potential overlap involving TCR sequences, wild-type YLQ, and individual 
mutant peptide sequences.

% INSERT EXACT LEAKAGE CRITERIA HERE.
% Possible categories:
% (1) identical mutant peptide in training data;
% (2) identical TCR in training data;
% (3) highly similar TCR according to the selected CDR3 similarity criterion;
% (4) prior exposure to wild-type YLQ.

Importantly, prior exposure to the wild-type YLQ antigen was distinguished 
from exposure to individual mutant peptides. Knowledge of the wild-type 
interaction may constitute a legitimate starting point for evaluating a 
model's ability to predict the effect of antigen mutation, whereas direct 
exposure to a specific mutant interaction may result in information leakage. 
We therefore recorded these forms of prior exposure separately and applied 
the predefined leakage-control criteria before calculating mutation-aware 
performance.

% INSERT EXACT FILTERING RULES AND NUMBERS REMOVED FOR EACH MODEL.



\subsection{Model Inference and Prediction Score Processing}

Each model was evaluated using its publicly available implementation and
recommended inference procedure. TCR sequences, mutant peptide sequences,
and additional inputs required by individual models were formatted according
to their original specifications.

Because the models produce prediction scores with different scales and may
use different model-specific decision thresholds, we evaluated predictive
performance using the partial area under the receiver operating characteristic
curve (pAUC). Specifically, pAUC was computed over the low-false-positive-rate
region with $\mathrm{FPR} \leq 0.1$. This threshold-independent evaluation
avoids reliance on model-specific classification cutoffs and provides a
consistent metric for comparing predictive discrimination across models and
experimental settings.

For model $m$, let $S_{ij}^{(m)}$ denote the prediction score for TCR $i$
and peptide $j$. The partial AUC was computed as

\begin{equation}
\mathrm{pAUC}_{0.1}^{(m)}
=
\int_{0}^{0.1}
\mathrm{TPR}_m(u)\,du,
\label{eq:model_pauc}
\end{equation}

where $\mathrm{TPR}_m(u)$ denotes the true-positive rate of model $m$ at
false-positive rate $u$. To facilitate comparison across models and settings,
we report the standardized partial AUC, for which a value of 0.5 corresponds
to random discrimination and 1.0 corresponds to perfect discrimination
within the specified FPR range.

Because prediction scores from different models may have different scales
and interpretations, raw scores were not directly compared across models.
For analyses of mutation-induced changes, we additionally considered
within-model score changes relative to the corresponding wild-type peptide:

\begin{equation}
\Delta S_{ij}^{(m)}
=
S^{(m)}(T_i,P_j)
-
S^{(m)}(T_i,P_{\mathrm{WT}}),
\label{eq:model_delta}
\end{equation}

where $P_{\mathrm{WT}}$ denotes the wild-type YLQPRTFLL peptide.
% INSERT:
% - software/model versions
% - handling of alpha/beta chains
% - handling of MHC inputs
% - model-specific preprocessing
% - models for which no published threshold exists

\subsection{Mutation-Aware Evaluation}

The central objective of the evaluation was to determine whether computational
TCR--epitope predictors reproduce experimentally observed patterns of
recognition following single-amino-acid antigen mutations. Rather than
evaluating each mutant independently, we organized the analysis at three
biologically meaningful levels: antigen position, experimentally defined
position groups, and individual TCR recognition fingerprints.

To provide a consistent measure of binary predictive performance across
models and experimental settings, we used the standardized partial area under
the receiver operating characteristic curve (pAUC) over the false-positive-rate
range $\mathrm{FPR} \leq 0.1$. The pAUC was computed directly from continuous
model prediction scores and experimental binary binding labels, avoiding the
use of model-specific classification thresholds. A standardized
$\mathrm{pAUC}_{0.1}$ of 0.5 corresponds to random discrimination, whereas
1.0 represents perfect discrimination within the specified FPR range.


\subsubsection{Position-level mutation analysis}

The primary mutation-aware analysis examined whether prediction models capture
the differential importance of individual positions in the nine-amino-acid
YLQ peptide. Mutant peptides were grouped according to the position of the
substituted residue, yielding nine position-specific sets corresponding to
P1--P9.

For each position $k \in \{1,\ldots,9\}$, all eligible TCR--mutant
interactions involving substitutions at that position were pooled, and
predictive performance was quantified using standardized
$\mathrm{pAUC}_{0.1}$:

\begin{equation}
\mathrm{pAUC}_{0.1}^{(k)}
=
\mathrm{pAUC}_{0.1}
\left(
\{y_{ij}: \operatorname{pos}(P_j)=k\},
\{S_{ij}: \operatorname{pos}(P_j)=k\}
\right),
\label{eq:position_pauc}
\end{equation}

where $y_{ij}$ is the experimentally observed binary binding label,
$S_{ij}$ is the corresponding model prediction score, and
$\operatorname{pos}(P_j)$ denotes the position altered in mutant peptide
$P_j$.

In parallel, the experimentally observed and model-predicted effects of
substitutions were summarized at each position relative to the wild-type
peptide. These position-level profiles were used to determine whether models
identify the peptide positions at which substitutions have the greatest
effects on TCR recognition. In particular, we examined whether model
predictions reproduce the experimentally observed contrast between relatively
mutation-tolerant positions and positions at which substitutions strongly
disrupt recognition.


\subsubsection{Grouped positional analysis}

To provide a more interpretable summary of positional effects, positions were
additionally organized according to the experimentally observed recognition
patterns reported by Malone et al.~\cite{malone2025fingerprinting}. Positions
were grouped into broadly mutation-tolerant, intermediate or TCR-dependent,
and highly restrictive categories based on the experimental characterization,
without using these categories during model inference.

For each position group, all eligible interactions associated with mutations
at positions in that group were pooled, and predictive performance was
quantified using the same standardized $\mathrm{pAUC}_{0.1}$ metric. We
additionally compared the relative mutation effects predicted for the three
groups with their experimentally observed ordering. A model that captures the
biological structure of the mutation landscape should assign greater
disruption to substitutions at restrictive positions than to substitutions at
tolerant positions. For example, mutations at P5 should generally produce a
larger predicted loss of recognition than mutations at the relatively tolerant positions if the model reproduces the experimentally observed
positional constraints.


\subsubsection{TCR-specific recognition fingerprints}
Among the groups, because the group of positions 4/6/7/8 are TCR-dependent, average group behavior is not enough.
We therefore performed a complementary TCR-specific analysis using the complete
set of eligible antigen variants for each receptor.

For TCR $i$, the experimentally measured responses across mutant peptides
define its recognition fingerprint. The corresponding computational
fingerprint was constructed from the model-predicted scores for the same
TCR--peptide pairs. Agreement between the experimental and predicted
fingerprints was quantified using Spearman's rank correlation coefficient:

\begin{equation}
\rho_i
=
\operatorname{Spearman}
\left(
\{B_{ij}\}_{j=1}^{M_i},
\{S_{ij}\}_{j=1}^{M_i}
\right),
\label{eq:fingerprint_spearman}
\end{equation}

where $M_i$ denotes the number of eligible peptide variants for TCR $i$,
$B_{ij}$ is the experimental recognition measurement, and $S_{ij}$ is the
corresponding model prediction score.

This analysis tests whether a model reproduces the relative pattern of
mutation sensitivity for individual receptors rather than only the
population-average positional effects. It is particularly relevant to
TCR-dependent positions such as P4, P6, P7, and P8, for which averaging
across receptors may mask distinct recognition patterns. Experimental and
predicted fingerprints were additionally visualized as mutation heatmaps.


\subsection{Statistical Analysis}

Standardized partial ROC-AUC with $\mathrm{FPR} \leq 0.1$ was used as the
primary binary prediction metric throughout the global, position-level, and
grouped analyses. The metric was calculated from continuous prediction scores
and therefore did not require model-specific classification thresholds.
Spearman's rank correlation coefficient was used separately for the
TCR-specific fingerprint analysis because this analysis evaluates agreement
in the relative pattern of mutation responses rather than binary
discrimination.

Performance statistics were calculated using [INSERT SOFTWARE/PACKAGE AND
VERSION]. Confidence intervals for $\mathrm{pAUC}_{0.1}$ and Spearman
correlations were estimated using [INSERT BOOTSTRAP PROCEDURE].

For comparisons involving multiple antigen positions, position groups, TCRs,
or prediction models, statistical significance was assessed using
[INSERT STATISTICAL TESTS]. Where multiple hypotheses were evaluated
simultaneously, $P$-values were adjusted using [INSERT METHOD, e.g.,
Benjamini--Hochberg false discovery rate correction]. A two-sided adjusted
$P$-value below [INSERT THRESHOLD] was considered statistically significant.

% INSERT:
% - number of bootstrap iterations
% - random seed if applicable
% - Python/R versions
% - package versions
% - code repository and data availability
\section{Results}

\end{document}